\documentclass[10pt,a4paper]{amsart}
\usepackage[margin=19mm]{geometry}
\usepackage{amsmath,amssymb,mathtools}
\usepackage[scr=rsfs,cal=euler]{mathalfa}
\usepackage{xfrac}
\usepackage[unicode,hidelinks,bookmarks=false]{hyperref}
\newcommand{\ie}{i.e.\ }
\newcommand{\cf}{cf.\ }
\newcommand{\resp}{respectively\ }
\newcommand{\Subsection}{\subsection}
\providecommand{\nobreakdash}{\nobreak\hbox{-}}
\setlength{\emergencystretch}{2em}
\multlinegap0pt
\usepackage{tikz}
\usepackage{amssymb}
\usepackage{xcolor}
\newtheorem{lem}{Lemma}
\newtheorem{thm}{Theorem}
\newtheorem{pf}{Pf}
\title{{The second and third Hankel determinants for certain classes of functions}}
\author{Vasudevarao Allu, Shobhit Kumar}
\date{Source: arXiv:2604.10301v1 (CC BY 4.0)}
\begin{document}
\maketitle

\noindent\emph{Source and licence.} {The second and third Hankel determinants for certain classes of functions}. Version arXiv:2604.10301v1 (CC BY 4.0), distributed by arXiv under the Creative Commons Attribution 4.0 International licence, http://creativecommons.org/licenses/by/4.0/, as recorded in the arXiv licence metadata for this exact version and shown on the abstract page https://arxiv.org/abs/2604.10301v1. Full licence notice: \texttt{licenses/arxiv-2604.10301-CC-BY-4.0.txt}.\par\smallskip

\noindent\textit{Editorial scope.} Counted unit: the article's thm thm (source0.tex lines 837--839). The article's complete original proof of it is delivered with it (source0.tex lines 840--1344, one \texttt{proof} environment of 505 lines). No mathematics is written, repaired or re-derived: the statement and the proof are the article's own lines, in the article's own notation, and no retained line is substituted or re-flowed.  Retained uncounted as the source-local context the proof needs: lines 554--561; lines 812--834.  Nothing was omitted from any retained span, so the package carries no omission record.  No retained line contains a citation, so the wrapper carries no bibliography and no bibliography entry.  No retained line contains a figure environment, an image-inclusion command or drawing code, so no image is delivered and the package declares no asset dependency.  Editorial formatting only, in addition: an AMS \texttt{amsart} preamble that restates the article's own declarations and macros used by the retained lines, namely the environments \texttt{lem}, \texttt{thm}, \texttt{pf} on separate counters, together with the standard packages those lines need.  The retained environments print in this document as Lemma 1, Theorem 1, Pf 1 in order of appearance, whereas the article prints the same items with the numbering of its own full document.  The complete native source of the article as distributed in the arXiv e-print for this exact version is bundled unchanged as \texttt{sources/198260/source0.tex}.
\begin{equation}\label{eq:3.6}
1+\frac{z f''(z)}{f'(z)}
=
1+\omega(z)+\frac{1}{4}\omega(z)^2
=
\varphi(\omega(z)),
\qquad z\in \mathbb{D}.
\end{equation}

\begin{lem}\label{lemma3}
Let
$
\omega(z)=\sum_{n=1}^{\infty}c_n z^n
$
be a Schwarz function, that is, $\omega$ is analytic in $\mathbb D$,
$\omega(0)=0$, and $|\omega(z)|<1$ for all $z\in\mathbb D$.
If $c_1\ge 0$, then there exist complex numbers $\gamma,\eta,\rho$
with
$
|\gamma|\le 1,\, |\eta|\le 1,\, |\rho|\le 1,
$
such that
\begin{align*}
c_2 &= (1-c_1^2)\gamma,\\[2mm]
c_3 &= (1-c_1^2)\bigl(\eta(1-|\gamma|^2)-c_1\gamma^2\bigr),\\[2mm]
c_4 &= (1-c_1^2)\Bigl(
c_1^2\gamma^3
-(1-|\gamma|^2)\bigl(2c_1\gamma\eta+\overline{\gamma}\eta^2\bigr)
+(1-|\gamma|^2)(1-|\eta|^2)\rho
\Bigr).
\end{align*}
\end{lem}

\begin{thm}
Let $f\in \mathcal{C}(\varphi)$. Then $|H_3(1)|\leq {1}/{144}$, where $H_3(1)$ denotes the third Hankel determinant. The result is sharp.
\end{thm}

\begin{pf}
Let $f\in \mathcal{C}(\varphi)$ be given by 
\[
1+z\frac{f''(z)}{f'(z)} \prec \varphi(z), \quad z\in \mathbb{D}.
\]
The third Hankel determinant is given  by
\begin{align}\label{eq:3.18}
H_3(1) =
\begin{vmatrix}
	1 & a_2 & a_3 \\
	a_2 & a_3 & a_4 \\
	a_3 & a_4 & a_5
\end{vmatrix}
= a_3(a_2 a_4 - a_3^2) - a_4(a_4 - a_2 a_3) + a_5(a_3 - a_2^2).
\end{align}
For the third Hankel determinant we use the coefficients of the Schwarz function together with the relation between the Carath\'eodory and Schwarz functions:
\[
\omega(z):=\frac{p(z)-1}{p(z)+1}, \quad z\in \mathbb{D},
\]
where
\[
\omega(z)=\sum_{n=1}^{\infty}c_n z^n
= c_1 z+c_2 z^2+c_3 z^3+\cdots.
\]
Now comparing the coefficients in \eqref{eq:3.6}, while using  $\omega(z)$ in the series form as above, we obtain
\begin{align}\label{eq:3.19}
\left.
\begin{aligned}
	a_{2} & = \frac{c_{1}}{2},\\[1mm]
	a_{3} & = \frac{1}{24}\left(5c_{1}^{2}+4c_{2}\right),\\[1mm]
	a_{4} & = \frac{1}{24}\left(-2c_{1}^{3}+c_{1}c_{2}+\frac{3}{4}c_{1}\left(5c_{1}^{2}+4c_{2}\right)+2c_{3}\right),\\[1mm]
	a_{5} & = \frac{43c_{1}^{4}+184c_{1}^{2}c_{2}+72c_{2}^{2}+176c_{1}c_{3}+96c_{4}}{1920}.
\end{aligned}
\right\}
\end{align}
Now putting \eqref{eq:3.19} in \eqref{eq:3.18}, we obtain
\begin{align}\label{eq:3.20}
69120 H_{3}(1)&=-7c_{1}^{6}+42c_{1}^{4}c_{2}+96c_{1}^{3}c_{3}+96c_{1}c_{2}c_{3}
-12c_{1}^{2}\left(17c_{2}^{2}+12c_{4}\right)\notag\\[2mm]
&\quad+16\left(7c_{2}^{3}-30c_{3}^{2}+36c_{2}c_{4}\right).
\end{align}
Since the class is rotation invariant, we may take $c_{1} \geq 0$. Moreover, since \(\omega\) is a Schwarz function, the Schwarz lemma gives \(|c_1|\le 1\). Hence
\[
0\le c_1\le 1.
\]
Now we use the parametric relation between the coefficients of the Schwarz functions in Lemma~\ref{lemma3} to obtain
\begin{align}
69120\,  H_{3}(1)&=-7c_{1}^{6}
+42\gamma\,c_{1}^{4}(1-c_{1}^{2})+96c_{1}^{3}\bigl(\eta(1-|\gamma|^{2})-\gamma^{2}c_{1}\bigr)(1-c_{1}^{2})\notag\\
&
\quad+96\gamma\,c_{1}\bigl(\eta(1-|\gamma|^{2})-\gamma^{2}c_{1}\bigr)(1-c_{1}^{2})^{2} \notag\\
&\quad
-12c_{1}^{2}\Bigl(
17\gamma^{2}(1-c_{1}^{2})^{2}
+12(1-c_{1}^{2})\notag\\
&\qquad
\Bigl(\gamma^{3}c_{1}^{2}
-(1-|\gamma|^{2})\bigl(-\rho(1-|\eta|^{2})+2\gamma\eta c_{1}+\eta^{2}\overline{\gamma}\bigr)\Bigr)
\Bigr) \notag\\
&\quad
+16\Bigl(
-30\bigl(\eta(1-|\gamma|^{2})-\gamma^{2}c_{1}\bigr)^{2}(1-c_{1}^{2})^{2}
+7\gamma^{3}(1-c_{1}^{2})^{3} \notag\\
&\quad
+36\gamma(1-c_{1}^{2})^{2}\Bigl(\gamma^{3}c_{1}^{2}
-(1-|\gamma|^{2})\bigl(-\rho(1-|\eta|^{2})+2\gamma\eta c_{1}+\eta^{2}\overline{\gamma}\bigr)\Bigr)
\Bigr).
\end{align}
We may re-write the above equation in the form
\begin{align}\label{eq:3.22}
69120\,  H_3(1)=A_{1}+B_{1} \eta+C_{1} \eta^2+D_{1} \rho,
\qquad  \mbox{where} \quad |\gamma|\le 1,\ |\eta|\le 1,\ |\rho|\le 1.
\end{align}
where
\begin{align*}
A_{1}&=-7c_{1}^{6}-42\gamma\,c_{1}^{4}(-1+c_{1}^{2})
+96\gamma^{4}c_{1}^{2}(-1+c_{1}^{2})^{2}
-12\gamma^{2}c_{1}^{2}\left(17-26c_{1}^{2}+9c_{1}^{4}\right)\\
&\quad-16\gamma^{3}\left(-7+27c_{1}^{2}-24c_{1}^{4}+4c_{1}^{6}\right),\\[2mm]
B_{1}&=96\left(-1+|\gamma|^{2}\right)c_{1}(-1+c_{1}^{2})
\left(\gamma+c_{1}^{2}+2\gamma c_{1}^{2}+2\gamma^{2}(-1+c_{1}^{2})\right),\\[2mm]
C_{1}&=-48\left(-1+|\gamma|^{2}\right)(-1+c_{1}^{2})
\Bigl(-10(-1+c_{1}^{2})+10|\gamma|^{2}(-1+c_{1}^{2})
-3\bigl(c_{1}^{2}+4\gamma(-1+c_{1}^{2})\bigr)\overline{\gamma}\Bigr),\\[2mm]
D_{1}&=144\left(-1+|\gamma|^{2}\right)\left(-1+|\eta|^{2}\right)(-1+c_{1}^{2})
\left(c_{1}^{2}+4\gamma(-1+c_{1}^{2})\right).
\end{align*}
Taking modulus on both sides in \eqref{eq:3.22}, writing
\[
p_1:=c_1,\qquad x:=|\gamma|\in[0,1],\qquad y:=|\eta|\in[0,1],
\]
and using $|\rho|\le 1$ together with the triangle inequality, we obtain
\begin{align*}
69120\, |H_{3}(1)| \leq H(p_1, x, y),
\end{align*}
where
\begin{align}\label{eq:3.23}
H(p_1,x,y)
&= 7p_1^{6}
+ x\Bigl(42p_1^{4}(1-p_1^{2})\Bigr)
+ x^{2}\Bigl(12p_1^{2}\bigl(17-26p_1^{2}+9p_1^{4}\bigr)\Bigr)\notag\\
&\quad
+ x^{3}\Bigl(16\bigl(7(1-p_1^{2})+20p_1^{2}(1-p_1^{2})+4p_1^{4}(1-p_1^{2})\bigr)\Bigr)
+ x^{4}\Bigl(96p_1^{2}(-1+p_1^{2})^{2}\Bigr)\notag\\
&\quad
+ y\Bigl(96(1-x^{2})p_1(1-p_1^{2})\bigl(x+p_1^{2}+2xp_1^{2}+2x^{2}(1-p_1^{2})\bigr)\Bigr)\notag\\
&\quad
+ y^{2}\Bigl(48(1-x^{2})(1-p_1^{2})\bigl(10(1-x^{2})(1-p_1^{2})
+3\bigl(p_1^{2}+4x(1-p_1^{2})\bigr)x\bigr)\Bigr)\notag\\
&\quad
+ 144(1-x^{2})(1-y^{2})(1-p_1^{2})\bigl(p_1^{2}+4x(1-p_1^{2})\bigr).
\end{align}
over the cuboid $\mathcal{D}:=\{(p_1, x, y):(p_1, x, y)\in [0, 1] \times [0, 1] \times [0, 1]\}$, which we denote by $\mathcal{D}$. Our aim is to maximize the expression $H(p_1, x, y)$ over this region $\mathcal{D}$.\\[2mm]
It is easy to see that the coefficient of $y$ remains non-negative in $\mathcal{D}$. Therefore, we take $y=1$ in the term
\[
y\Bigl(96(1-x^{2})p_1(1-p_1^{2})\bigl(x+p_1^{2}+2xp_1^{2}+2x^{2}(1-p_1^{2})\bigr)\Bigr)
\]
to make the computation easier. Furthermore, we define a new function $H_{1}(p_1, x, y)$ such that $H(p_1, x, y) \leq H_{1}(p_1, x, y)$ in $\mathcal{D}$. Precisely,
\begin{align}\label{eq:3.24}
H_1(p_1,x,y)
&= 7p_1^{6}
+ x\Bigl(42p_1^{4}(1-p_1^{2})\Bigr)
+ x^{2}\Bigl(12p_1^{2}\bigl(17-26p_1^{2}+9p_1^{4}\bigr)\Bigr)\notag\\
&\quad
+ x^{3}\Bigl(16\bigl(7(1-p_1^{2})+20p_1^{2}(1-p_1^{2})+4p_1^{4}(1-p_1^{2})\bigr)\Bigr)
+ x^{4}\Bigl(96p_1^{2}(-1+p_1^{2})^{2}\Bigr)\notag\\
&\quad
+ 96(1-x^{2})p_1(1-p_1^{2})\bigl(x+p_1^{2}+2xp_1^{2}+2x^{2}(1-p_1^{2})\bigr)\notag\\
&\quad
+ y^{2}\Bigl(48(1-x^{2})(1-p_1^{2})\bigl(10(1-x^{2})(1-p_1^{2})
+3\bigl(p_1^{2}+4x(1-p_1^{2})\bigr)x\bigr)\Bigr)\notag\\
&\quad
+ 144(1-x^{2})(1-y^{2})(1-p_1^{2})\bigl(p_1^{2}+4x(1-p_1^{2})\bigr).
\end{align}
Since $H_{1}(p_{1}, x, y)$ is quadratic in $y$ and the coefficient of $y$ is $0$, the maximum occurs at one of the endpoints $y=0$ or $y=1$. We handle these two cases separately.\\[2mm]
Set $G_{1}(p_1, x):=H_{1}(p_{1}, x, 1)$ and $G_{2}(p_1, x):=H_{1}(p_{1}, x, 0)$. \\[2mm]
For the case when $y=1$ we prove that the function $G_{1}(p_1, x)\le 480$ in $[0, 1]\times[0, 1]$. For the same purpose we take $480-G_{1}(p_{1}, x)$ and prove that it remains non-negative in $[0, 1] \times[0, 1]$.\\
\[
\begin{aligned}
G_1(p_1,x)=\;&7 p_1^6 + 42 p_1^4 (1-p_1^2)x + 12 p_1^2(17-26p_1^2+9p_1^4)x^2 \\
&+16\bigl(7(1-p_1^2)+20p_1^2(1-p_1^2)+4p_1^4(1-p_1^2)\bigr)x^3 \\
&+96p_1^2(-1+p_1^2)^2x^4 \\
&+96p_1(1-p_1^2)(1-x^2)\bigl(p_1^2+x+2p_1^2x+2(1-p_1^2)x^2\bigr) \\
&+48(1-p_1^2)(1-x^2)\bigl(3x(p_1^2+4(1-p_1^2)x)+10(1-p_1^2)(1-x^2)\bigr).
\end{aligned}
\]
Taking
\[
p:=p_1,
\qquad
F(p,x):=480-G_1(p,x),
\]
we show that
\[
F(p,x)\ge 0
\qquad \text{for all }(p,x)\in[0,1]\times[0, 1],
\]
with equality only at $(0,0)$.
We use the fact that if a polynomial $H(u,v)$ is written on the unit square in the  Bernstein basis
\[
H(u,v)=\sum_{i=0}^{m}\sum_{j=0}^{n} c_{ij}\,B_i^m(u)B_j^n(v),
\]
then
\[
\min_{i,j} c_{ij}\le H(u,v)\le \max_{i,j} c_{ij}
\qquad 0\le u,v\le 1,
\]
because all Bernstein basis functions are nonnegative and their sum is $1$.
Since $G_1$ has bidegree $(6,4)$, we use the  Bernstein basis of bidegree $(6,4)$.\\[2mm]
First we divide the square \([0, 1]\times[0, 1]\) into
\begin{align*}
Q_1&=\left[0,\frac12\right]\times\left[0,\frac12\right],\qquad
Q_2=\left[0,\frac12\right]\times\left[\frac12,1\right],\\[3mm]
Q_3&=\left[\frac12,1\right]\times\left[0,\frac12\right],\qquad
Q_4=\left[\frac12,1\right]\times\left[\frac12,1\right].
\end{align*}
%%%%%%%%%%
For each subrectangle \(Q_k=[a,b]\times[c,d]\), we first transfer the problem to the unit square by writing
\[
p=a+(b-a)u,\qquad x=c+(d-c)v,\qquad (u,v)\in[0, 1]\times[0, 1].
\]
The polynomial \(G_1(p,x)\) then becomes a polynomial in \((u,v)\), which we express in the Bernstein basis of bidegree \((6,4)\),
\[
G_1(a+(b-a)u,\;c+(d-c)v)
=
\sum_{i=0}^{6}\sum_{j=0}^{4} b_{ij}\,B_i^6(u)B_j^4(v).
\]
Because
\[
B_i^6(u)\ge 0,\qquad B_j^4(v)\ge 0,
\qquad
\sum_{i=0}^{6}\sum_{j=0}^{4} B_i^6(u)B_j^4(v)=1,
\]
it follows that the value of the polynomial on the rectangle is bounded above by the largest Bernstein coefficient \(b_{ij}\) arising from that representation. Performing this calculation on each of the four rectangles gives
\begin{align*}
G_1(p,x) &\le 481 \qquad \text{on } Q_1,\\[2mm]
G_1(p,x) &\le 398 \qquad \text{on } Q_2,\\[2mm]
G_1(p,x) &\le \frac{75535}{256} \quad \text{on } Q_3,\\[2mm]
G_1(p,x) &\le \frac{18463}{64} \quad \text{on } Q_4.
\end{align*}
%%%%%%%%%%
Since
\begin{align*}
398<480,\qquad \frac{75535}{256}<480,\qquad \frac{18463}{64}<480,
\end{align*}
it follows that
\[
F(p,x)>0
\qquad \text{on }Q_2\cup Q_3\cup Q_4.
\]
Thus only $Q_{1}$ can contain a zero of $F$.
Now subdivide
\[
Q_1=\left[0,\frac12\right]\times\left[0,\frac12\right]
\]
into
\begin{align*}
Q_{11}&=\left[0,\frac14\right]\times\left[0,\frac14\right],\qquad
Q_{12}=\left[0,\frac14\right]\times\left[\frac14,\frac12\right],\\[3mm]
Q_{13}&=\left[\frac14,\frac12\right]\times\left[0,\frac14\right],\qquad
Q_{14}=\left[\frac14,\frac12\right]\times\left[\frac14,\frac12\right].
\end{align*}
Applying the Bernstein method once more on these four rectangles, we obtain
\begin{align*}
G_1(p,x) &\le \frac{1921}{4} \qquad \text{on } Q_{11},\\[2mm]
G_1(p,x) &\le \frac{14681}{32} \qquad \text{on } Q_{12},\\[2mm]
G_1(p,x) &\le \frac{13939571}{32768} \quad \text{on } Q_{13},\\[2mm]
G_1(p,x) &\le \frac{13594541}{32768} \quad \text{on } Q_{14}.
\end{align*}
Since
\begin{align*}
\frac{14681}{32}<480,\qquad
\frac{13939571}{32768}<480,\qquad
\frac{13594541}{32768}<480,
\end{align*}
we get
\[
F(p,x)>0
\qquad \text{on }Q_{12}\cup Q_{13}\cup Q_{14}.
\]
Therefore it remains only to show that
\[
F(p,x)\ge 0
\qquad \text{on }Q_{11}=\left[0,\frac14\right]\times\left[0,\frac14\right].
\]
We now consider the rectangle $Q_{11}$.
First expand $F(p_1, x)$ by 
\[
\begin{aligned}
F(p,x)=\;&960p^2-96px+384x^2 -96p^3-480p^4+96p^5-7p^6 \\
&+\bigl(-144p^2-96p^3+102p^4+192p^5+42p^6\bigr)x \\
&+\bigl(-192p-972p^2+480p^3+696p^4-288p^5-108p^6\bigr)x^2 \\
&+\bigl(-112+96p-64p^2+96p^3+112p^4-192p^5+64p^6\bigr)x^3 \\
&+\bigl(96+192p-288p^2-384p^3+288p^4+192p^5-96p^6\bigr)x^4.
\end{aligned}
\]
Let, 
\[
F(p,x)=Q(p,x)+R(p,x),
\]
where,
\[
Q(p,x)=960p^2-96px+384x^2
\]
and
\[
\begin{aligned}
R(p,x)=\;&-96p^3-480p^4+96p^5-7p^6 \\
&+\bigl(-144p^2-96p^3+102p^4+192p^5+42p^6\bigr)x \\
&+\bigl(-192p-972p^2+480p^3+696p^4-288p^5-108p^6\bigr)x^2 \\
&+\bigl(-112+96p-64p^2+96p^3+112p^4-192p^5+64p^6\bigr)x^3 \\
&+\bigl(96+192p-288p^2-384p^3+288p^4+192p^5-96p^6\bigr)x^4.
\end{aligned}
\]
Every monomial in $R(p, x)$ has total degree at least $3$.
We first estimate $Q(p, x)$ from below. Since
$
2px\le p^2+x^2,
$
we have
$
96px\le 48p^2+48x^2.
$
Therefore
\[
Q(p,x)\ge (960-48)p^2+(384-48)x^2
=912p^2+336x^2
\ge 336(p^2+x^2).
\]
Next we estimate $R(p, x)$ on $Q_{11}$, where
\[
0\le p\le \frac14,\qquad 0\le x\le \frac14.
\]
Let
\begin{align*}
R(p,x)=\sum_{i+j\ge 3} c_{ij}p^i x^j.
\end{align*}
For every monomial $p^i x^j$ with $i+j\ge 3$, we claim that
\begin{align*}
p^i x^j \le 4^{-(i+j-2)}(p^2+x^2).
\end{align*}
Indeed, if \(i\ge 2\), then
\begin{align*}
p^i x^j=p^2 p^{\,i-2}x^j\le 4^{-(i+j-2)}p^2
\le 4^{-(i+j-2)}(p^2+x^2).
\end{align*}
If $i=1$, then $j\ge 2$, and
\begin{align*}
px^j=x^2(px^{j-2})\le 4^{-(j-1)}x^2
=4^{-(i+j-2)}x^2
\le 4^{-(i+j-2)}(p^2+x^2).
\end{align*}
If \(i=0\), then \(j\ge 3\), and
\[
x^j\le 4^{-(j-2)}x^2
=4^{-(i+j-2)}x^2
\le 4^{-(i+j-2)}(p^2+x^2).
\]
Therefore, we have 
\begin{align*}
|R(p,x)|
\le
\left(\sum_{i+j\ge 3}|c_{ij}|\,4^{-(i+j-2)}\right)(p^2+x^2).
\end{align*}
Now group the coefficients according to the total degree \(i+j\).
The sums of the absolute values of the coefficients of total degree
\(3,4,5,6,7,8,9,10\) are respectively
\[
544,\quad 1740,\quad 934,\quad 1279,\quad 826,\quad 588,\quad 256,\quad 96.
\]
Hence
\[
\sum_{i+j\ge 3}|c_{ij}|\,4^{-(i+j-2)}
=
\frac{544}{4}
+\frac{1740}{16}
+\frac{934}{64}
+\frac{1279}{256}
+\frac{826}{1024}
+\frac{588}{4096}
+\frac{256}{16384}
+\frac{96}{65536}.
\]
A direct simplification gives
\[
\sum_{i+j\ge 3}|c_{ij}|\,4^{-(i+j-2)}
=
\frac{543349}{2048}.
\]
Therefore
\[
|R(p,x)|\le \frac{543349}{2048}(p^2+x^2).
\]
Combining the estimates for $Q$ and $R$, we get
\[
F(p,x)=Q(p,x)+R(p,x)
\ge 336(p^2+x^2)-\frac{543349}{2048}(p^2+x^2).
\]
Thus
\[
F(p,x)\ge
\left(336-\frac{543349}{2048}\right)(p^2+x^2)
=
\frac{144779}{2048}(p^2+x^2).
\]
Since
\[
\frac{144779}{2048}>0,
\]
it follows that
\[
F(p,x)\ge 0
\qquad \text{for all }(p,x)\in Q_{11},
\]
it is easy to see that
\begin{align*}
F(p,x)>0
\qquad \text{for }(p,x)\in Q_{11}\setminus\{(0,0)\}.
\end{align*}
Also, we have
\[
F(0,0)=480-G_1(0,0)=0.
\]
Combining this with the Bernstein bounds on the other rectangles, we conclude that
\[
F(p,x)\ge 0
\qquad \text{for all }(p,x)\in[0, 1]\times[0, 1],
\]
with equality only at \((0,0)\).
Equivalently,
\[
G_1(p_1,x)\le 480
\qquad \text{for all }(p_1,x)\in[0, 1]\times[0, 1],
\]
and equality holds only at
\[
(p_1,x)=(0,0).
\]
Therefore
\[
\max_{0\le p_1\le 1,\;0\le x\le 1}G_1(p_1,x)=480.
\]
Now we solve the case when 
$y=0$, precisely for the function $G_{2}(p_1, x):=H_{1}(p_1, x, 0)$.\\[2mm]
Setting $p_{1}=p$, and applying the Bernstein method to
\[
\begin{aligned}
G_2(p,x)=\;&7 p^6 + 42 p^4 (1-p^2)x + 12 p^2(17-26p^2+9p^4)x^2 \\
&+16\bigl(7(1-p^2)+20p^2(1-p^2)+4p^4(1-p^2)\bigr)x^3 \\
&+96p^2(-1+p^2)^2x^4 \\
&+144(1-p^2)\bigl(p^2+4(1-p^2)x\bigr)(1-x^2) \\
&+96p(1-p^2)(1-x^2)\bigl(p^2+x+2p^2x+2(1-p^2)x^2\bigr)
\end{aligned}
\]
on the rectangle
\[
[0,1]\times[0,1].
\]
Since the given rectangle is already the unit square, no affine transformation is needed. We first expand \(G_2(p,x)\) in the power basis:
\[
\begin{aligned}
G_2(p,x)=\;&\bigl(7p^6-96p^5-144p^4+96p^3+144p^2\bigr) \\
&+\bigl(-42p^6-192p^5+618p^4+96p^3-1152p^2+96p+576\bigr)x \\
&+\bigl(108p^6+288p^5-168p^4-480p^3+60p^2+192p\bigr)x^2 \\
&+\bigl(-64p^6+192p^5-832p^4-96p^3+1360p^2-96p-464\bigr)x^3 \\
&+\bigl(96p^6-192p^5-192p^4+384p^3+96p^2-192p\bigr)x^4.
\end{aligned}
\]
Now consider the Bernstein basis of bidegree \((6,4)\):
\[
B_i^6(p)=\binom{6}{i}p^i(1-p)^{6-i},
\qquad
B_j^4(x)=\binom{4}{j}x^j(1-x)^{4-j}.
\]
Then
\[
G_2(p,x)=\sum_{i=0}^{6}\sum_{j=0}^{4} c_{ij}\,B_i^6(p)B_j^4(x),
\]
where the Bernstein coefficients \(c_{ij}\) are determined by the power-basis expansion.

For the present polynomial, the Bernstein coefficient matrix \((c_{ij})\) is \\
\[
\begin{pmatrix}
0 & 144 & 288 & 316 & 112 \\[5mm]
0 & 148 & \dfrac{904}{3} & 340 & 112 \\[5mm]
\dfrac{48}{5} & \dfrac{712}{5} & \dfrac{4298}{15} & \dfrac{1022}{3} & \dfrac{2188}{15} \\[5mm]
\dfrac{168}{5} & \dfrac{666}{5} & \dfrac{1234}{5} & \dfrac{1566}{5} & \dfrac{1068}{5} \\[5mm]
\dfrac{336}{5} & \dfrac{1271}{10} & \dfrac{2917}{15} & \dfrac{7639}{30} & \dfrac{802}{3} \\[5mm]
80 & \dfrac{215}{2} & 127 & \dfrac{935}{6} & \dfrac{634}{3} \\[5mm]
7 & 7 & 7 & 7 & 7
\end{pmatrix}.
\]\\
We now use the fact that for \((p,x)\in[0, 1]\times[0, 1]\),
\[
B_i^6(p)\ge 0,\qquad B_j^4(x)\ge 0,
\]
and
\[
\sum_{i=0}^{6}\sum_{j=0}^{4} B_i^6(p)B_j^4(x)=1.
\]
Hence \(G_2(p,x)\) is a convex combination of the Bernstein coefficients. Therefore,
\[
\min_{0\le i\le 6,\;0\le j\le 4} c_{ij}
\le
G_2(p,x)
\le
\max_{0\le i\le 6,\;0\le j\le 4} c_{ij}
\]
for all $(p,x)\in[0, 1]\times[0, 1]$.
From the coefficient matrix above, we see that
\[
\max_{0\le i\le 6,\;0\le j\le 4} c_{ij}=\frac{1022}{3}.
\]
Consequently,
\[
G_2(p_1,x)\le \frac{1022}{3}
\qquad
\text{for all }0\le p_1\le 1,\ 0\le x\le 1.
\]
Thus the Bernstein method gives the upper bound
\[
\max_{0\le p_1\le 1,\;0\le x\le 1} G_2(p_1,x)\le \frac{1022}{3}\le 480.
\]
Hence, exhausting all the cases, we have
\[
69120\,|H_3(1)|\le 480.
\]
Therefore,
\begin{align*}
|H_3(1)|\le \frac{480}{69120}=\frac{1}{144}.
\end{align*}
Finally, we show that $\omega(z)=z^3$ provides an extremal function for the third Hankel determinant. In this case,
\[
1+z\frac{\widetilde{f}''(z)}{\widetilde{f}'(z)}=\varphi(z^3)
=1+z^3+\frac{z^6}{4}.
\]
Therefore,
\begin{align*}
\widetilde{f}(z)
&=\int_{0}^{z} \exp\!\left(\int_{0}^{\zeta} \frac{\varphi(t^3)-1}{t}\,dt\right)\, d\zeta,\\[2pt]
\widetilde{f}(z)
&=z + \frac{1}{12}z^4 + \frac{1}{72} z^7+\ldots .
\end{align*}
It is easy to check that this $\widetilde{f}$ gives $|H_3(1)|=\dfrac{1}{144}$, which proves the sharpness.
\end{pf}

\end{document}
